\documentclass[11pt]{article}

\usepackage[utf8]{inputenc}
\usepackage[T1]{fontenc}
\usepackage{lmodern}
\usepackage{geometry}
\usepackage{amsmath,amssymb,amsfonts,amsthm,mathtools}
\usepackage{bm}
\usepackage{enumitem}
\usepackage{microtype}
\usepackage{hyperref}
\usepackage{titlesec}
\usepackage{booktabs}
\usepackage{array}
\usepackage{setspace}
\usepackage{mathrsfs}
\usepackage{physics}

\geometry{margin=1in}
\hypersetup{colorlinks=true, linkcolor=black, urlcolor=blue, citecolor=black}
\setlist[enumerate]{topsep=0.4em,itemsep=0.35em}
\setlist[itemize]{topsep=0.3em,itemsep=0.25em}
\titleformat{\section}{\large\bfseries}{\thesection.}{0.5em}{}
\titleformat{\subsection}{\normalsize\bfseries}{\thesubsection.}{0.5em}{}
\setstretch{1.06}

\newcommand{\Aether}{\AE{}ther}
\newcommand{\Aflow}{\AE{}ther-flow}
\newcommand{\TheoryName}{\Aflow{} Interpretation of Relativity}
\newcommand{\TheoryProgram}{\Aether{} / \Aflow{} framework}

\newcommand{\CompletedMetric}{g^{\sharp}}
\newcommand{\MatterSpace}{\Psi}

\newcommand{\Car}{\mathrm{car}}
\newcommand{\Cur}{\mathrm{cur}}
\newcommand{\Hom}{\mathrm{hom}}

\newcommand{\ForSpace}[1]{\mathcal I_{#1}^{\mathrm{for}}}
\newcommand{\PrimShell}{\mathcal S_{\mathrm{prim}}}
\newcommand{\Reduction}[1]{\mathcal R_{#1}}
\newcommand{\CurrentCarrier}[1]{\mathfrak C_{#1}^{\mathrm{cur}}}
\newcommand{\EqCarrier}[1]{\mathfrak C_{#1}^{\mathrm{eq}}}
\newcommand{\MetricOut}{\mathscr G_{\mathrm{out}}}
\newcommand{\MatterOut}{\mathscr T_{\mathrm{out}}}

\newcommand{\FoundSpace}[1]{\mathcal F_{#1}^{\mathrm{fdn}}}
\newcommand{\GroundSpace}[1]{\mathcal G_{#1}^{\mathrm{grd}}}
\newcommand{\OrigSpace}[1]{\mathcal O_{#1}^{\mathrm{orig}}}
\newcommand{\PrecSpace}[1]{\mathcal P_{#1}^{\mathrm{prec}}}
\newcommand{\LiftSpace}[1]{\mathcal L_{#1}^{\mathrm{lift}}}
\newcommand{\SubSpace}[1]{\mathcal M_{#1}^{\mathrm{sub}}}
\newcommand{\SubTarget}[1]{E_{#1}^{\mathrm{sub}}}
\newcommand{\SubPair}[1]{\mathfrak s_{#1}^{\mathrm{sub}}}
\newcommand{\EqnTarget}[1]{Q_{#1}^{\mathrm{eqn}}}

\newcommand{\HiddenSpace}[1]{\mathcal H_{#1}^{\mathrm{hid}}}
\newcommand{\HiddenBody}[1]{D_{#1}^{\mathrm{hid}}}

\newcommand{\SeedPackage}{\mathfrak S^{\mathrm{seed}}}
\newcommand{\SeedSpace}[1]{\mathcal S_{#1}^{\mathrm{seed}}}
\newcommand{\SeedBody}[1]{\mathcal B_{#1}^{\mathrm{seed}}}
\newcommand{\SeedProj}[1]{\Pi_{#1}^{\mathrm{seed}}}
\newcommand{\SeedFoundMin}[1]{\mathfrak f_{#1}^{\mathrm{seed,fdn}}}
\newcommand{\SeedGroundMin}[1]{\mathfrak f_{#1}^{\mathrm{seed,grd}}}
\newcommand{\SeedOrigMin}[1]{\mathfrak f_{#1}^{\mathrm{seed,orig}}}
\newcommand{\SeedPrecMin}[1]{\mathfrak f_{#1}^{\mathrm{seed,prec}}}
\newcommand{\SeedLiftMin}[1]{\mathfrak f_{#1}^{\mathrm{seed,lift}}}
\newcommand{\SeedSubMin}[1]{\mathfrak m_{#1}^{\mathrm{seed}}}
\newcommand{\SeedSection}[1]{\Gamma_{#1}^{\mathrm{seed}}}
\newcommand{\SeedFunctional}[1]{\mathscr W_{#1}^{\mathrm{seed}}}
\newcommand{\SeedShellSet}[1]{\mathcal Z_{#1}^{\mathrm{seed}}}
\newcommand{\SeedZeroCore}[1]{\mathcal Z_{#1}^{0,\mathrm{seed}}}
\newcommand{\SeedSplit}[1]{\Phi_{#1}^{\mathrm{seed}}}
\newcommand{\SeedCharge}[1]{\mathcal Q_{#1}^{\mathrm{seed}}}
\newcommand{\SeedEmbed}[1]{\zeta_{#1}^{\mathrm{seed}}}
\newcommand{\SeedKernelCh}[1]{\widetilde K_{#1}^{0,\mathrm{seed}}}
\newcommand{\SeedStateComp}[1]{P_{#1}^{\mathrm{seed,co}}}

\newcommand{\GermPackage}{\mathfrak G^{\mathrm{germ}}}
\newcommand{\GermSpace}[1]{\mathcal G_{#1}^{\mathrm{germ}}}
\newcommand{\GermBody}[1]{\mathcal B_{#1}^{\mathrm{germ}}}
\newcommand{\GermProj}[1]{\Pi_{#1}^{\mathrm{germ}}}
\newcommand{\GermSeedMin}[1]{\mathfrak f_{#1}^{\mathrm{germ,seed}}}
\newcommand{\GermFoundMin}[1]{\mathfrak f_{#1}^{\mathrm{germ,fdn}}}
\newcommand{\GermGroundMin}[1]{\mathfrak f_{#1}^{\mathrm{germ,grd}}}
\newcommand{\GermOrigMin}[1]{\mathfrak f_{#1}^{\mathrm{germ,orig}}}
\newcommand{\GermPrecMin}[1]{\mathfrak f_{#1}^{\mathrm{germ,prec}}}
\newcommand{\GermLiftMin}[1]{\mathfrak f_{#1}^{\mathrm{germ,lift}}}
\newcommand{\GermSubMin}[1]{\mathfrak m_{#1}^{\mathrm{germ}}}
\newcommand{\GermSection}[1]{\Gamma_{#1}^{\mathrm{germ}}}
\newcommand{\GermFunctional}[1]{\mathscr W_{#1}^{\mathrm{germ}}}
\newcommand{\GermShellSet}[1]{\mathcal Z_{#1}^{\mathrm{germ}}}
\newcommand{\GermZeroCore}[1]{\mathcal Z_{#1}^{0,\mathrm{germ}}}
\newcommand{\GermSplit}[1]{\Phi_{#1}^{\mathrm{germ}}}
\newcommand{\GermCharge}[1]{\mathcal Q_{#1}^{\mathrm{germ}}}
\newcommand{\GermEmbed}[1]{\zeta_{#1}^{\mathrm{germ}}}
\newcommand{\GermKernelCh}[1]{\widetilde K_{#1}^{0,\mathrm{germ}}}
\newcommand{\GermStateComp}[1]{P_{#1}^{\mathrm{germ,co}}}

\newcommand{\RootNucleus}[1]{\mathfrak N_{#1}^{\mathrm{rt}}}

\newtheorem{definition}{Definition}
\newtheorem{proposition}{Proposition}
\newtheorem{remark}{Remark}
\newtheorem{corollary}{Corollary}
\newtheorem{theorem}{Theorem}

\title{\textbf{The \TheoryName{}}\\[0.4em]
\large Primitive-Reservoir Observer-Reduction\\
\large Nucleus-Generating Substrate Seed Dynamics\\
\large from Seed-Generating Substrate Germ Dynamics}
\author{}
\date{}

\begin{document}

\maketitle

\begin{abstract}
The current primitive-reservoir observer-reduction frontier is the theorem deriving the minimal root exact-shell transport nucleus from nucleus-generating substrate seed dynamics. The accompanying seed-frontier routing note then fixed the next honest move: either derive or justify the seed package from still more primitive substrate structure, or prove a genuinely smaller theorem-level collapse strictly inside the observer-relevant seed package. The present note takes the default first route.

For each sector it introduces \emph{seed-generating substrate germ dynamics}: a compact convex germ body over the recorded seed state space, unique germ minimizers on the fixed seed, foundation, ground, origin, precursor, lift, and substrate fibers, a smooth germ restriction section, an exact germ shell projecting diffeomorphically to the exact seed shell, a closed embedded germ zero core whose projected split reproduces the seed split, a hidden charge, an observer-compatible exact germ zero-response realization, and a fixed-data solvability / coercive package projecting to the recorded seed kernel data. The current nucleus-generating substrate seed dynamics are then recovered canonically as the projected exact-shell transport image of that deeper germ package.

The benchmark boundary remains fixed throughout. The one operative metric is still exactly \(\CompletedMetric\), the ordinary matter route is unchanged, there is no second operative metric, and the low-energy matter law is not enlarged. The positive result is therefore narrow but benchmark facing: the current seed package is no longer left as an unexplained primitive stopping point on the active line. The remaining honest burden moves one layer deeper again, to the germ dynamics themselves or to a sharper collapse inside their observer-relevant core.
\end{abstract}

\section{Introduction}

The active derivational program of \TheoryName{} has now reached a sharper burden than another same-output seed restatement. The current seed theorem already showed that the minimal root exact-shell transport nucleus is the exact-shell transport image of a more primitive nucleus-generating substrate seed dynamics package. The accompanying routing note then fixed the consequence of that theorem: the next honest benchmark-facing manuscript must either derive or justify the seed package from still more primitive substrate structure, or prove a genuinely smaller collapse strictly inside the observer-relevant seed package. Reopening the old root-nucleus handoff, re-running another same-output root / foundation / ground relay, or returning to stationary-reduction, stationary-Hessian, spectral, or selector layers would no longer address the live burden.

The present note takes the first route. Its positive claim is not that final substrate microphysics has been identified, and it is not that the benchmark exact-relativistic package has been changed. Its claim is narrower and benchmark facing. The current seed package is shown to arise canonically from a more primitive germ-level substrate dynamics whose projected exact-shell transport already contains the observer-relevant structure of the seed package. In that sense the seed layer is no longer the primitive place where the active line has to stop.

Two features matter here. First, the theorem targets the current seed package itself rather than revisiting the already-spent root-nucleus handoff, so it does not spend the next move on an older same-output presentation. Second, the deeper germ package is not merely another renaming of the seed package. It adds an explicit germ state space, germ body, germ-to-seed projection, and deeper exact-shell transport data below the seed frontier. The seed shell, seed zero core, split, and lower minimizer hierarchy are therefore projected images of a deeper substrate package rather than unexplained primitive data.

\paragraph{Framework and claim boundary.}
\TheoryProgram{} denotes the broader \Aether{} / \Aflow{} framework built on the claims that the \Aether{} is the underlying four-dimensional substrate of reality, the \Aflow{} is its intrinsic ordered motion, observed three-dimensional space is the local experiential slice of that deeper substrate, S-time is the experienced order of change arising from matter, light, and the \Aflow{}, and observed expansion is the three-dimensional appearance of deeper four-dimensional ordered motion. Within that framework, \TheoryName{} denotes the exact-closure theory in which the effective gravitational dynamics are adopted to be exactly Einsteinian with universal matter coupling. In the active sequence, \emph{adoption} means use of that established relativistic dynamics without claiming substrate derivation, while \emph{derivation} is reserved for a first-principles recovery from explicit substrate variables.

The present note stays strictly inside that boundary. It does not alter the benchmark exact-closure package. It does not introduce a second operative metric or a new low-energy matter law. It only proves that the current nucleus-generating substrate seed dynamics can itself be generated from a more primitive germ-level substrate dynamics.

\section{Fixed Inputs}

\begin{definition}[Recorded primitive continuation data]
\label{def:fixed}
Fix the following continuation data from the active primitive-reservoir observer-reduction line.
\begin{enumerate}
    \item For each sector label \(\sigma\in\{\Car,\Cur,\Hom\}\), a primitive forcing shell
    \[
    \ForSpace{\sigma},
    \]
    a visible reduction map
    \[
    \Reduction{\sigma}:\ForSpace{\sigma}\to X_{\sigma},
    \]
    and shell-level current and equilibrium carriers
    \[
    \CurrentCarrier{\sigma}:\ForSpace{\sigma}\to Z_{\sigma}^{\mathrm{cur}},
    \qquad
    \EqCarrier{\sigma}:\ForSpace{\sigma}\to Z_{\sigma}^{\mathrm{eq}}.
    \]
    \item The product shell
    \[
    \PrimShell
    \equiv
    \ForSpace{\Car}\times \ForSpace{\Cur}\times \ForSpace{\Hom}.
    \]
    \item The recorded metric and ordinary-matter outputs
    \[
    \MetricOut:\PrimShell\to\{\text{observer metrics}\},
    \qquad
    \MatterOut:\PrimShell\times\MatterSpace\to\{\text{observer stress tensors}\},
    \]
    satisfying
    \begin{equation}
    \MetricOut(\mathbf w)=\CompletedMetric,
    \qquad
    \MatterOut(\mathbf w,\psi)
    =
    T_{\mu\nu}^{\mathrm{matter}}[\CompletedMetric,\psi]
    \label{eq:fixedoutputs}
    \end{equation}
    for every \(\mathbf w\in\PrimShell\) and every matter configuration \(\psi\in\MatterSpace\).
\end{enumerate}
\end{definition}

\begin{definition}[Recorded lower-fiber ladder]
\label{def:ladder}
Fix \(\sigma\in\{\Car,\Cur,\Hom\}\). The active line already fixes the lower-fiber ladder
\[
\FoundSpace{\sigma},
\qquad
\GroundSpace{\sigma},
\qquad
\OrigSpace{\sigma},
\qquad
\PrecSpace{\sigma},
\qquad
\LiftSpace{\sigma},
\qquad
\SubSpace{\sigma},
\]
together with the common substrate target \(\SubTarget{\sigma}\), the positive-definite pairing
\[
\SubPair{\sigma}:
\SubTarget{\sigma}\times\SubTarget{\sigma}\to\mathbb R,
\]
and the hidden equation target \(\EqnTarget{\sigma}\). The current seed frontier uses exactly this recorded ladder when it names the seed projection and the six minimizer images that remain benchmark facing at seed scope.
\end{definition}

\begin{remark}
Definitions \ref{def:fixed} and \ref{def:ladder} are not reopened here. The primitive continuation data, the one operative metric, the ordinary matter route, and the recorded lower-fiber ladder remain fixed inputs. The present note addresses only whether the current nucleus-generating substrate seed dynamics can itself be generated from more primitive substrate dynamics.
\end{remark}

\section{Recorded Seed Target}

\begin{definition}[Recorded nucleus-generating substrate seed dynamics]
\label{def:seed}
Fix \(\sigma\in\{\Car,\Cur,\Hom\}\). The recorded seed package \(\SeedPackage\) for sector \(\sigma\) consists of the following data.
\begin{enumerate}
    \item A finite-dimensional Euclidean seed state space \(\SeedSpace{\sigma}\), a compact convex seed body
    \[
    \SeedBody{\sigma}\subset\SeedSpace{\sigma}
    \]
    with nonempty interior, an affine surjection
    \[
    \SeedProj{\sigma}:\SeedSpace{\sigma}\to\FoundSpace{\sigma},
    \]
    and smooth seed fiber minimizers
    \[
    \SeedFoundMin{\sigma},\quad
    \SeedGroundMin{\sigma},\quad
    \SeedOrigMin{\sigma},\quad
    \SeedPrecMin{\sigma},\quad
    \SeedLiftMin{\sigma},\quad
    \SeedSubMin{\sigma}
    \]
    on the fixed foundation, ground, origin, precursor, lift, and substrate fibers of the recorded ladder.
    \item A smooth seed restriction section
    \[
    \SeedSection{\sigma}:\SeedSpace{\sigma}\to\SubTarget{\sigma},
    \]
    the seed functional
    \[
    \SeedFunctional{\sigma}(s)
    \equiv
    \SubPair{\sigma}\bigl(\SeedSection{\sigma}(s),\SeedSection{\sigma}(s)\bigr),
    \]
    and the exact seed shell
    \[
    \SeedShellSet{\sigma}
    \equiv
    \bigl(\SeedSection{\sigma}\bigr)^{-1}(0)\cap\SeedBody{\sigma}.
    \]
    \item A closed embedded seed zero core
    \[
    \SeedZeroCore{\sigma}\subset\SeedShellSet{\sigma},
    \]
    a finite-dimensional hidden fiber space \(\HiddenSpace{\sigma}\) with compact hidden body \(\HiddenBody{\sigma}\subset\HiddenSpace{\sigma}\), and a split diffeomorphism
    \[
    \SeedSplit{\sigma}:
    \SeedZeroCore{\sigma}\times\HiddenBody{\sigma}
    \xrightarrow{\cong}
    \SeedShellSet{\sigma}.
    \]
    \item A hidden charge
    \[
    \SeedCharge{\sigma}:\HiddenSpace{\sigma}\to\EqnTarget{\sigma},
    \]
    a smooth observer-compatible exact seed zero-response realization
    \[
    \SeedEmbed{\sigma}:\ForSpace{\sigma}\hookrightarrow\SeedZeroCore{\sigma},
    \]
    and the fixed-data solvability / coercive package on \(\SeedZeroCore{\sigma}\), recorded through the charted kernel \(\SeedKernelCh{\sigma}\) and orthogonal projector \(\SeedStateComp{\sigma}\).
\end{enumerate}
\end{definition}

\begin{remark}
Definition \ref{def:seed} is the fixed target already carried by the current seed-frontier theorem. The present note does not spend its move on a second seed-level restatement. It records the current seed package only because the deeper germ dynamics are defined relative to it.
\end{remark}

\section{Seed-Generating Substrate Germ Dynamics}

\begin{definition}[Seed-generating substrate germ dynamics]
\label{def:germ}
Fix \(\sigma\in\{\Car,\Cur,\Hom\}\). A seed-generating substrate germ dynamics package \(\GermPackage\) for sector \(\sigma\) consists of the following data.
\begin{enumerate}
    \item A finite-dimensional Euclidean germ state space \(\GermSpace{\sigma}\), a compact convex germ body
    \[
    \GermBody{\sigma}\subset\GermSpace{\sigma}
    \]
    with nonempty interior, an affine surjection
    \[
    \GermProj{\sigma}:\GermSpace{\sigma}\to\SeedSpace{\sigma},
    \]
    and smooth germ fiber minimizers
    \[
    \GermSeedMin{\sigma},\quad
    \GermFoundMin{\sigma},\quad
    \GermGroundMin{\sigma},\quad
    \GermOrigMin{\sigma},\quad
    \GermPrecMin{\sigma},\quad
    \GermLiftMin{\sigma},\quad
    \GermSubMin{\sigma}
    \]
    satisfying
    \[
    \GermProj{\sigma}\bigl(\GermBody{\sigma}\bigr)=\SeedBody{\sigma},
    \qquad
    \GermProj{\sigma}\circ\GermSeedMin{\sigma}
    =
    \mathrm{id}_{\SeedBody{\sigma}},
    \]
    together with the transport identities
    \[
    \GermProj{\sigma}\circ\GermFoundMin{\sigma}
    =
    \SeedFoundMin{\sigma},
    \qquad
    \GermProj{\sigma}\circ\GermGroundMin{\sigma}
    =
    \SeedGroundMin{\sigma},
    \]
    \[
    \GermProj{\sigma}\circ\GermOrigMin{\sigma}
    =
    \SeedOrigMin{\sigma},
    \qquad
    \GermProj{\sigma}\circ\GermPrecMin{\sigma}
    =
    \SeedPrecMin{\sigma},
    \]
    \[
    \GermProj{\sigma}\circ\GermLiftMin{\sigma}
    =
    \SeedLiftMin{\sigma},
    \qquad
    \GermProj{\sigma}\circ\GermSubMin{\sigma}
    =
    \SeedSubMin{\sigma}.
    \]
    \item A smooth germ restriction section
    \[
    \GermSection{\sigma}:\GermSpace{\sigma}\to\SubTarget{\sigma},
    \]
    the germ functional
    \[
    \GermFunctional{\sigma}(g)
    \equiv
    \SubPair{\sigma}\bigl(\GermSection{\sigma}(g),\GermSection{\sigma}(g)\bigr),
    \]
    and the exact germ shell
    \[
    \GermShellSet{\sigma}
    \equiv
    \bigl(\GermSection{\sigma}\bigr)^{-1}(0)\cap\GermBody{\sigma},
    \]
    satisfying
    \[
    \GermSection{\sigma}\circ\GermSeedMin{\sigma}
    =
    \SeedSection{\sigma},
    \qquad
    \GermProj{\sigma}\big|_{\GermShellSet{\sigma}}
    :
    \GermShellSet{\sigma}
    \xrightarrow{\cong}
    \SeedShellSet{\sigma}.
    \]
    \item A closed embedded germ zero core
    \[
    \GermZeroCore{\sigma}\subset\GermShellSet{\sigma},
    \]
    the same hidden fiber space \(\HiddenSpace{\sigma}\) with the same compact hidden body \(\HiddenBody{\sigma}\subset\HiddenSpace{\sigma}\), a diffeomorphism
    \[
    \GermProj{\sigma}\big|_{\GermZeroCore{\sigma}}
    :
    \GermZeroCore{\sigma}
    \xrightarrow{\cong}
    \SeedZeroCore{\sigma},
    \]
    and a split diffeomorphism
    \[
    \GermSplit{\sigma}:
    \GermZeroCore{\sigma}\times\HiddenBody{\sigma}
    \xrightarrow{\cong}
    \GermShellSet{\sigma}
    \]
    obeying the projected split identity
    \[
    \SeedSplit{\sigma}(z,\eta)
    =
    \GermProj{\sigma}\Bigl(
    \GermSplit{\sigma}\bigl(
    (\GermProj{\sigma}\big|_{\GermZeroCore{\sigma}})^{-1}(z),\eta
    \bigr)\Bigr)
    \]
    for every \(z\in\SeedZeroCore{\sigma}\) and every \(\eta\in\HiddenBody{\sigma}\).
    \item A hidden charge
    \[
    \GermCharge{\sigma}:\HiddenSpace{\sigma}\to\EqnTarget{\sigma},
    \]
    a smooth observer-compatible exact germ zero-response realization
    \[
    \GermEmbed{\sigma}:\ForSpace{\sigma}\hookrightarrow\GermZeroCore{\sigma},
    \]
    satisfying
    \[
    \SeedCharge{\sigma}
    =
    \GermCharge{\sigma},
    \qquad
    \SeedEmbed{\sigma}
    =
    \GermProj{\sigma}\circ\GermEmbed{\sigma},
    \]
    and a fixed-data solvability / coercive package on \(\GermZeroCore{\sigma}\), recorded through \(\GermKernelCh{\sigma}\) and \(\GermStateComp{\sigma}\), whose projection through
    \[
    \GermProj{\sigma}\big|_{\GermZeroCore{\sigma}}
    \]
    recovers the recorded seed kernel chart \(\SeedKernelCh{\sigma}\) and orthogonal projector \(\SeedStateComp{\sigma}\).
\end{enumerate}
The package is \emph{benchmark faithful} if it preserves the benchmark outputs \eqref{eq:fixedoutputs}, introduces no second operative metric, introduces no enlarged low-energy matter law, and is presented as a deeper substrate-side source for the current seed package rather than as a replacement benchmark theory.
\end{definition}

\begin{remark}
The label \emph{germ} is used here only as a local marker for a deeper layer below the current seed frontier. It does not by itself assert a completed microscopic ontology or a final primitive law. Its role is only to name the next sharper transport package below the current seed dynamics.
\end{remark}

\section{Projected Exact-Shell Transport to the Seed Package}

\begin{proposition}[The seed shell is recovered by projected exact-shell transport]
\label{prop:shell}
Fix a benchmark-faithful germ package. Then projected exact-shell transport along
\[
\GermProj{\sigma}\big|_{\GermShellSet{\sigma}}
\qquad\text{and}\qquad
\GermProj{\sigma}\big|_{\GermZeroCore{\sigma}}
\]
recovers the exact seed shell, the embedded seed zero core, and the seed split of Definition \ref{def:seed}.
\end{proposition}

\begin{proof}
Definition \ref{def:germ}(2) states that the restriction of \(\GermProj{\sigma}\) to \(\GermShellSet{\sigma}\) is a diffeomorphism onto \(\SeedShellSet{\sigma}\). Definition \ref{def:germ}(3) states that the restriction of \(\GermProj{\sigma}\) to \(\GermZeroCore{\sigma}\) is a diffeomorphism onto \(\SeedZeroCore{\sigma}\) and that the projected image of the germ split is exactly the recorded seed split. These are precisely the claimed projected exact-shell transport statements.
\end{proof}

\begin{proposition}[Germ dynamics canonically reproduce the current seed package]
\label{prop:seed}
Fix a benchmark-faithful germ package. Then the projected exact-shell transport of Definition \ref{def:germ} reproduces exactly the current nucleus-generating substrate seed dynamics package of Definition \ref{def:seed}.
\end{proposition}

\begin{proof}
The seed body and the lower minimizer hierarchy are recovered by Definition \ref{def:germ}(1): the germ body projects onto the recorded seed body, the germ seed minimizer sections the projection on that body, and the projected germ lower minimizers are exactly the recorded seed minimizers on the fixed foundation, ground, origin, precursor, lift, and substrate fibers.

The seed restriction data are recovered by Definition \ref{def:germ}(2): the recorded seed section is the minimizing transport of the germ section, the seed functional is therefore the projected restriction functional, and Proposition \ref{prop:shell} recovers the exact seed shell itself.

The remaining seed shell-internal data are recovered by Definition \ref{def:germ}(3)--(4): the embedded seed zero core is the projected germ zero core, the seed split is the projected germ split, the seed hidden charge is the germ hidden charge, the exact seed zero-response realization is the projected germ realization, and the fixed-data solvability / coercive package is the projection of the corresponding germ package through the zero-core diffeomorphism. These are exactly the constituents retained by the current seed frontier.
\end{proof}

\begin{proposition}[The benchmark package remains fixed]
\label{prop:boundary}
For every benchmark-faithful germ package, the operative metric remains exactly \(\CompletedMetric\), the ordinary matter route remains unchanged, no second operative metric is introduced, and the low-energy matter law is not enlarged.
\end{proposition}

\begin{proof}
This is part of the benchmark-faithfulness requirement in Definition \ref{def:germ}. Equation \eqref{eq:fixedoutputs} remains unchanged on the full primitive forcing shell, so the deeper germ package does not alter the benchmark exact-closure package.
\end{proof}

\section{Main Theorem}

\begin{theorem}[Nucleus-generating substrate seed dynamics from seed-generating substrate germ dynamics]
\label{thm:main}
Fix the primitive continuation data of Definition \ref{def:fixed}, the recorded lower-fiber ladder of Definition \ref{def:ladder}, and the recorded seed package of Definition \ref{def:seed}. Then, for each sector \(\sigma\in\{\Car,\Cur,\Hom\}\), every benchmark-faithful seed-generating substrate germ dynamics package canonically produces the current seed package as projected exact-shell transport from the deeper germ layer. More precisely:
\begin{enumerate}
    \item the current nucleus-generating substrate seed dynamics are reproduced unchanged, sector by sector, by projected exact-shell transport from the deeper germ package;
    \item because the induced seed package is exactly the one already used by the current seed-frontier theorem, the minimal root exact-shell transport nucleus and the previously recorded downstream root-side consequences remain available unchanged;
    \item the same benchmark one-metric package remains fixed: the operative metric is still exactly \(\CompletedMetric\), the ordinary matter route is unchanged, there is no second operative metric, and the low-energy matter law is not enlarged;
    \item consequently the remaining honest burden on the active line is deeper again: derive or justify the seed-generating substrate germ dynamics themselves from still more primitive substrate structure, or else prove a genuinely smaller theorem-level collapse at that germ level.
\end{enumerate}
\end{theorem}

\begin{proof}
Item (1) is Proposition \ref{prop:seed}. Item (2) is the routing consequence of item (1): the present note reproduces exactly the seed package that the current seed-frontier theorem already uses as its input, so that theorem and its downstream recovery of the minimal root exact-shell transport nucleus apply unchanged. Item (3) is Proposition \ref{prop:boundary}. Item (4) is the project-level consequence of the previous items.
\end{proof}

\begin{corollary}[What remains open after this theorem]
\label{cor:next}
After Theorem \ref{thm:main}, the remaining honest burden is no longer to justify the nucleus-generating substrate seed dynamics themselves. Those dynamics are now the canonical projected exact-shell output of a sharper germ class. What remains open is why the seed-generating substrate germ dynamics should hold, or whether an even sharper theorem can remove that still deeper germ-level freedom while preserving the same benchmark one-metric package and claim boundary.
\end{corollary}

\begin{proof}
Theorem \ref{thm:main} removes the current seed package as the unexplained stopping point on the active line. The next honest burden therefore lies strictly below it, unless a smaller theorem-level core inside the germ package can first be isolated and proved sufficient.
\end{proof}

\section{Conclusion}

The positive result of this note is narrower than a completed first-principles derivation and stronger than another same-output seed restatement. The current seed-frontier theorem had already shown that the minimal root exact-shell transport nucleus is the exact-shell transport image of a more primitive seed package. The present theorem then removes the next unexplained stopping point by showing that this seed package itself is the projected exact-shell transport image of a more primitive seed-generating substrate germ dynamics.

The scope remains disciplined. The benchmark package is unchanged: the one operative metric is still exactly \(\CompletedMetric\), the ordinary matter route is unchanged, there is no second operative metric, and the low-energy matter law is not enlarged. Nothing here upgrades adoption to completed derivation or licenses stronger promotion language.

The open burden therefore moves one honest step further again. The next benchmark-facing manuscript should derive or justify the germ dynamics from still more primitive substrate structure, unless the sharper honest gain available instead is a genuinely smaller theorem-level collapse inside the observer-relevant germ package. Until then, the present result should be read for what it is: a deeper seed-scope derivation on the active line of \TheoryName{}, not a claim that the full first-principles program is finished.

\end{document}
